\section{System model and problem formulation}
\label{sec:system}

This section describes the near-RT RIC considered in this work and the logging requirement on which the proposed
approach rests. It then formulates the two problems that the paper addresses: learning a signed causal effect map from
logged decisions, and using that map to mitigate conflicts within safety limits.

\subsection{Near-RT RIC and conflict mitigation}
\label{sec:setting_ric}

The system consists of a near-RT RIC that hosts several xApps and controls a set of cells (Figure~\ref{fig:system_e6p}).
Each xApp reads measurement reports from the cells and sends control requests that change one of its control
parameters on a cell, and each family of control parameters has a single writer. Every request passes a
conflict-mitigation interface in the style of the O-RAN Working Group~3 report~\cite{ORANConflictMit}, which accepts or
rejects it, as the near-RT RIC platform does with a control request that an xApp sends over the E2
interface~\cite{ORANRICARCH}. We call the rule that makes this choice the referee. A rejected request is left
unapplied, and the xApp may propose it again later.
The network state is summarised by per-cell KPIs, such as protected-slice violations, service violations, energy,
radio-link failures (RLFs) and served load. A request on a cell acts on that cell, on its neighbours, taken from the
cell's observable neighbour list, and on the remaining cells; these are the own, neighbour and far relations.

\begin{figure}[tbp]
\centering
\includegraphics[width=\textwidth]{figures/fig_system_e6p.pdf}
\caption{System model.}
\label{fig:system_e6p}
\end{figure}

\subsection{Decision units and logging}
\label{sec:units}

The referee decides per decision unit rather than per request. A unit $u$ is identified by a cell, an xApp and an
opening time: it opens at the xApp's first request on that cell while no unit of the pair is active. Its mode
$W_u$, accept or reject, is drawn once at opening and holds for every request of the pair for 60~s. An accepted unit
applies the requests of its cell-xApp pair during this hold, and a rejected unit applies none. The referee accepts with
the propensity $\pi_u$ and rejects with probability $1 - \pi_u$, and the logging requirement of this paper is that the
RIC records, for every unit, its mode and its propensity. Each unit also carries an observation-only context: the
opening request, the cell's recent reports and applied configuration, and their aggregates over the neighbouring
cells. Discovery methods read only the logged units, their context and the per-cell KPI series; the simulator's
internal counters and the ground truth remain hidden from them.

Three logging policies are considered. The randomised policy accepts or rejects each unit with probability 0.5,
independently, from the start of an episode. The incumbent policy is confounded: its propensity, within
$[0.15, 0.85]$, depends on the request class and on the cell state, namely the physical resource block (PRB)
utilisation for energy-saving requests and the fraction of protected user equipments (UEs) below their throughput
floor for protected-share requests. Placebo logs draw and record modes in the same way but apply every request, so
that no logged mode has any effect.

\subsection{Problem formulation}
\label{sec:estimand}

A hypothesis $(f, r, k)$ concerns a parameter family $f$, a relation $r$ and a KPI $k$; in the scenario of
Section~\ref{sec:design}, four families, three relations and five KPIs give $m = 60$ hypotheses. Because the referee
randomises the mode of a unit rather than a parameter value, each hypothesis concerns a mode contrast held for 60~s.
Let $y_u$ be the outcome of unit $u$ for $(r, k)$, namely the KPI summed over the cells of $r$ in the 90~s after
opening minus the same sum in the 90~s before. Let $\mathcal{F}_{u-1}$ denote the history of the network, of every xApp and of the
modes of earlier units up to the opening time. The xApp chooses the size of the change and the direction
$\mathrm{sgn}_u$ of the opening request, which orients the contrast, so that a positive effect means that accepting a
request that raises the parameter raises the KPI.

\begin{definition}[Policy-averaged assignment effect]
\label{def:tau_pi}
Let $u$ be a unit of family $f$. Its assignment effect is
\begin{equation}
\tau_u^{\pi} = \mathbb{E}_{\pi}\bigl[y_u \mid \mathcal{F}_{u-1}, W_u = \text{accept}\bigr]
             - \mathbb{E}_{\pi}\bigl[y_u \mid \mathcal{F}_{u-1}, W_u = \text{reject}\bigr],
\label{eq:tau_u}
\end{equation}
where $\mathbb{E}_{\pi}$ averages over the network's future randomness and over the modes that the logging policy
$\pi$ draws for later units. The policy-averaged assignment effect $\tau_{\pi}$ of hypothesis $(f, r, k)$ is the
limit of $\bar\tau_{\pi} = N^{-1} \sum_u \mathrm{sgn}_u\, \tau_u^{\pi}$ over the $N$ units of $f$ as the number of
episodes grows.
\end{definition}

The target depends on the logging policy, which draws the modes of the other units in a unit's window, so randomised and
incumbent logs target different values of $\tau_{\pi}$. The knockout ground truth of Section~\ref{sec:gt} targets a third contrast, $\tau_{\mathrm{AA}}$, in which every other unit is accepted
(Appendix~\ref{app:estimands}). No assumption of non-interference between cells or units is needed, because the null
hypothesis concerns the whole assignment vector of a family. Under an alternative, the effect combines the direct
effect of a unit's mode on the cells of its relation, the spillover from neighbouring units' modes, and the carryover
through the policy, since a rejected request opens a new unit after the hold.

The paper addresses two problems. The first is discovery: from the logged units and KPI series alone, and without a
model of the network, decide for each hypothesis whether the sharp null $H_0(f, r, k)$ holds. The sharp null states that the
outcome a unit of $f$ would have at any cell and opening time is the same under any assignment of $f$'s modes. Each
rejected null is reported as an edge signed by the estimated direction of $\tau_{\pi}$. The resulting signed causal effect map should control the false discovery
rate (FDR) over the sharp nulls at a level $q$. The second problem is mitigation: given such a map, choose a referee
that minimises the protected-slice violations $V$ subject to two safety constraints. First, the referee must keep at least 90~\% of
the energy saving that the energy-saving xApps achieve on their own. Second, each of four guard KPIs, namely the service
violation rate, non-protected broadband violations, low-latency violations and RLFs, must stay within 1.10 times its
value under the static rule that accepts every request (accept-all).
Section~\ref{sec:method} addresses both problems, and Section~\ref{sec:design} specifies the scenario in which they are
evaluated.

